\SourceLeaf{258}%
auf folgende Weise. Durch Differentiation der ersten und der zweiten der zu Beginn dieses Artikels angegebenen Gleichungen erhalten wir
\[
\sin\psi\mathord{.}\frac{\mathrm{d}n}{\mathrm{d}q}+n\cos\psi\mathord{.}\frac{\mathrm{d}\psi}{\mathrm{d}q}+\sin\psi\mathord{.}\frac{\mathrm{d}\psi}{\mathrm{d}p}=0
\]
woraus zusammen mit
\[
n\cos\psi\mathord{.}\frac{\mathrm{d}\varphi}{\mathrm{d}q}+\sin\psi\mathord{.}\frac{\mathrm{d}\varphi}{\mathrm{d}p}=0
\]
folgt
\[
\frac{r\sin\psi}{n}\mathord{.}\frac{\mathrm{d}n}{\mathrm{d}q}+\frac{r\sin\psi}{n}\mathord{.}\frac{\mathrm{d}(\psi+\varphi)}{\mathrm{d}p}+r\cos\psi\mathord{.}\frac{\mathrm{d}(\psi+\varphi)}{\mathrm{d}q}=0
\]
Aus dieser Gleichung gewinnen wir mit der Methode der unbestimmten Koeffizienten leicht die Reihe für \(\psi+\varphi\), wenn wir berücksichtigen, dass ihr erstes Glied \(\tfrac{1}{2}\pi\) sein muss, sofern der Radius als Einheit genommen wird und \(2\pi\) den Kreisumfang bezeichnet,
\[
\begin{aligned}
\relax[6]\quad\psi+\varphi={}&\tfrac{1}{2}\pi-f^0pq-\tfrac{2}{3}f'ppq-(\tfrac{1}{2}f''-\tfrac{1}{6}f^0f^0)p^3q\quad\text{etc.}\\
&-g^0pqq-\tfrac{3}{4}g'ppqq\\
&-(h^0-\tfrac{1}{3}f^0f^0)pq^3
\end{aligned}
\]

Es scheint sich zu lohnen, auch den Flächeninhalt des Dreiecks \(ABD\) in eine Reihe zu entwickeln. Dazu dient die folgende Bedingungsgleichung, die sich leicht aus recht naheliegenden geometrischen Überlegungen herleiten lässt und in der \(S\) den gesuchten Flächeninhalt bezeichnet:
\[
\frac{r\sin\psi}{n}\mathord{.}\frac{\mathrm{d}S}{\mathrm{d}p}+r\cos\psi\mathord{.}\frac{\mathrm{d}S}{\mathrm{d}q}=\frac{r\sin\psi}{n}\mathord{.}\int n\,\mathrm{d}q
\]
wobei die Integration bei \(q=0\) beginnt. Mit der Methode der unbestimmten Koeffizienten erhalten wir daraus
\[
\begin{aligned}
\relax[7]\quad S={}&\tfrac{1}{2}pq-\tfrac{1}{12}f^0p^3q-\tfrac{1}{20}f'p^4q-(\tfrac{1}{30}f''-\tfrac{1}{60}f^0f^0)p^5q\quad\text{etc.}\\
&-\tfrac{1}{12}f^0pq^3-\tfrac{3}{40}g^0p^3qq-\tfrac{1}{20}g'p^4qq\\
&-\tfrac{7}{120}f'ppq^3-(\tfrac{1}{15}h^0+\tfrac{2}{45}f''+\tfrac{1}{60}f^0f^0)p^3q^3\\
&-\tfrac{1}{10}g^0pq^4-\tfrac{3}{40}g'ppq^4\\
&-(\tfrac{1}{10}h^0-\tfrac{1}{30}f^0f^0)pq^5
\end{aligned}
\]
\Article{25}
Die Formeln des vorhergehenden Artikels beziehen sich auf ein rechtwinkliges Dreieck aus kürzesten Linien. Wir gehen nun zum allgemeinen Fall über. Es sei \(C\) ein weiterer Punkt auf derselben kürzesten Linie
